\documentclass[
    openany,
    12pt,
    a4paper,
    twoside,
    english,
    brazil, sumario=tradicional
]{/home/lucas_galdino/Desktop/TCC/planning/abntex2/abntex2}

\usepackage{lmodern}			% Usa a fonte Latin Modern			
\usepackage[T1]{fontenc}		% Selecao de codigos de fonte.
\usepackage[utf8]{inputenc}		% Codificacao do documento (conversão automática dos acentos)
\usepackage{indentfirst}		% Indenta o primeiro parágrafo de cada seção.
\usepackage{color}				% Controle das cores
\usepackage{graphicx}			% Inclusão de gráficos
\usepackage{microtype} 			% para melhorias de justificação

\usepackage[brazilian,hyperpageref]{backref}	 % Paginas com as citações na bibl
\usepackage[alf]{/home/lucas_galdino/Desktop/TCC/planning/abntex2/abntex2cite}	% Citações padrão ABNT

% my own changes
\usepackage{/home/lucas_galdino/Desktop/TCC/planning/abntex2/my_changes}

\titulo{Avaliação N1}
% \titlehead{{\centering\includegraphics[width=6cm]{../assets/ufsc-logo.png}}}
\autor{Lucas Eduardo Galdino Silva}
\local{Brasil}
\data{Outubro, 2024}
\instituicao{%
  Universidade Federal de Santa Catarina -- UFSC
  \par
  Engenharia Mecânica
}
\tipotrabalho{Avaliação N1}
% O preambulo deve conter o tipo do trabalho, o objetivo, 
% o nome da instituição e a área de concentração 
\preambulo{Avaliação N1 para conclusão da matéria Planejamento de Trabalho de Curso\\
}%  \imprimirorientadorRotulo~\imprimirorientador}


% informações do PDF
\makeatletter
\hypersetup{
    %pagebackref=true,
		pdftitle={\@title}, 
		pdfauthor={\@author},
    pdfsubject={\imprimirpreambulo},
	  pdfcreator={\@author},
		pdfkeywords={abnt}{latex}{abntex}{abntex2}{trabalho acadêmico}, 
		bookmarksdepth=4
}
% let chapter start in any page
\@openrightfalse
\makeatother

\renewcommand{\cleardoublepage}{}

% Makes abntex2 number the chapters with roman numbers
\AtBeginDocument{\aliaspagestyle{chapter}{plain}}

% Makes abntex2 do not start a new page before each chapter
\let\oldchapter\chapter
\renewcommand{\chapter}{\clearpage\oldchapter}

\usepackage{titlesec}

% Redefine the \chapter command to not start on a new page
\titleformat{\chapter}[hang]
  {\normalfont\huge\bfseries} % Format for chapter title
  {\thechapter} % Chapter number
  {20pt} % Space between number and title
  {\huge} % Format for chapter title text

\titlespacing*{\chapter}{0pt}{0pt}{20pt} % Adjust spacing before and after



\begin{document}

\selectlanguage{brazil}

\imprimircapa

\imprimirfolhaderosto*

\pdfbookmark[0]{\contentsname}{toc}
\tableofcontents
\cleardoublepage

\cleardoublepage
\textual



\chapter{Introdução}

O Problema de Roteamento de Veículos (VRP) é um problema de otimização combinatória que envolve a determinação de rotas eficientes para um conjunto de veículos, 
com o objetivo de minimizar a distância total percorrida ou o tempo total gasto. 
Esse problema é amplamente estudado na literatura de pesquisa operacional e ciência da computação, 
e tem aplicações em diversas áreas, como logística, transporte, entrega de mercadorias, coleta de lixo, serviços de emergência e muito mais.\\

É um problema NP-difícil, 
o que significa que não há uma solução eficiente conhecida para encontrar a solução ótima em tempo polinomial. 

É um braço de pesquisa do Problema do Caixeiro Viajante (TSP-Traveling Salesman Problem), problema clássico de análise combinatória 
e otimização, que teve suas primeiras pesquisas com esse nome em 1930, mas que foi formalizado e popularizado em 1956 por George Dantzig, 
Richard Bellman e Fulkerson. 

Durante seu livro, COOK (In Pursuit Of The Travelling Salesman, \citeyear{pursuit_travelling_salesman}) relata que o TSP é ainda um braço da teoria dos grafos,
 que é um ramo da matemática responsável por estudar estruturas 
abstratas chamadas grafos, que consistem em um conjunto de vértices (pontos) e um conjunto de arestas (linhas) que conectam esses vértices. O TSP vem do 
estudo de ciclos hamiltonianos. Eles possuem este nome devido ao matemático irlandês William Rowan Hamilton, que inventou o jogo icosiano, que consiste em buscar 
a uma maneira de visitar todos os vértices de um grafo, sem repetir nenhum, e retornar ao vértice de origem.

O projeto consiste em desenvolver e avaliar uma solução computacional avançada para otimização de rotas de caminhões de entrega, 
com foco na aplicação de métodos de aceleração por GPU (Unidade de processamento gráfico). \\
O objetivo principal é explorar a eficiência de algoritmos de roteamento de veículos em ambientes de alta complexidade, envolvendo múltiplos 
fatores que tornam o problema mais realista e aplicável ao mundo 
real. Essa solução busca resolver uma variante do Problema de Roteamento de Veículos (VRP), chamada Problema de Roteamento de Veículos com frota variada, 
com entrega e retirada e janelas de tempo ou HFVRPPDTW (Heterogeneous Fleet 
Vehicle Routing Problem with Pickup and Delivery and Time Windows), que apresenta restrições adicionais, como a diversidade das frotas, 
a existência de janelas de tempo para entrega e a necessidade de gerenciar tanto coleta quanto entrega em diferentes locais.

Essa versão específica do problema do VRP exige uma abordagem robusta para minimizar a distância total percorrida pela frota 
e, simultaneamente, maximizar o número de pedidos entregues dentro do prazo. Em um cenário de transporte logístico, o uso de 
frota heterogênea implica que cada veículo tem características e capacidades distintas, demandando um algoritmo que possa determinar 
não apenas quais rotas são mais curtas, mas também quais veículos devem ser alocados para cada trajeto, considerando tanto as restrições 
de tempo quanto a localização dos pontos de coleta e entrega.

Para explorar as vantagens da aceleração por GPU, a linguagem principal será \texttt{Python}, utilizando bibliotecas como \texttt{Numba} e 
\texttt{PyCUDA}, que facilitam a execução de operações intensivas em paralelo na GPU. O projeto também poderá integrar implementações em \texttt{C++} 
e \texttt{CUDA}, a fim de comparar o desempenho e a escalabilidade das soluções desenvolvidas em diferentes linguagens e abordagens de programação 
paralela. O uso dessas tecnologias permite que o projeto lide com a elevada demanda computacional imposta por grandes volumes de dados e múltiplas 
variáveis envolvidas na roteirização otimizada de veículos.

O armazenamento e processamento de dados poderão ser suportados por bancos de dados como \texttt{Dask} e \texttt{SQLite3}, selecionados por suas capacidades
de lidar com grandes volumes de dados em tempo real. Esses bancos de dados também poderão fornecer suporte para Sistemas de Informação Geográfica (GIS), 
essenciais para processar as coordenadas geoespaciais dos pontos de entrega e coleta, bem como para construir mapas e rotas que permitam uma visualização 
clara dos resultados.

A escolha pelo uso de softwares de código aberto (utilizando software livres sempre que possível) é outro aspecto importante do projeto. Além de contribuir para a democratização das tecnologias, 
permitindo que outros pesquisadores e desenvolvedores repliquem e ampliem os resultados, essa abordagem promove a inovação e a colaboração, 
possibilitando o acesso e a modificação do código para adaptação a diferentes cenários e necessidades específicas.


\chapter{Motivação}

% Fazer um for loop para uma função latex que gere lista de siglas automaticamente
A relevância do problema reside em sua aplicabilidade prática em diversos setores, como transporte, logística, distribuição de mercadorias e serviços,
podendo, a depender do setor e nível tecnológico de uma empresa, realizar uma diferença enorme. Além disso, o problema é um desafio computacional
significativo, com múltiplas restrições e variáveis que tornam a solução ótima uma tarefa complexa.

Sua ampla aplicabilidade prática em diversos setores, como transporte, logística, 
distribuição de mercadorias e prestação de serviços. Esses setores são pilares de operação para inúmeras empresas, e a 
otimização de rotas de entrega pode trazer melhorias significativas em termos de custos, tempo de entrega e sustentabilidade 
ambiental. Por exemplo, em uma cadeia de suprimentos moderna, onde as empresas competem para reduzir prazos de entrega e minimizar 
despesas, uma solução de roteamento eficiente pode reduzir o número de veículos necessários, otimizar o consumo de combustível e, 
consequentemente, diminuir as emissões de carbono. Dependendo do setor específico e do nível de maturidade tecnológica de uma empresa, 
esse tipo de solução pode fazer uma enorme diferença em termos de competitividade, qualidade de serviço e sustentabilidade. Em cenários 
onde há diversas restrições, como horários de entrega rígidos, rotas limitadas e veículos com capacidades diferentes, a otimização de 
rotas passa a ser um diferencial estratégico.

Além disso, o problema do roteamento de veículos é um desafio computacional significativo. Trata-se de um problema de 
otimização com múltiplas variáveis e restrições que tornam a obtenção de uma solução ótima uma tarefa complexa e demorada. 
Este problema pertence à classe dos problemas NP-difíceis, o que significa que o tempo de cálculo aumenta exponencialmente 
com o número de locais e veículos envolvidos. As variáveis são inúmeras: horários de atendimento em cada ponto, capacidades 
dos veículos, diferentes tipos de carga e até fatores externos, como o tráfego e condições climáticas, que precisam ser 
considerados para que a solução final seja eficaz. Assim, o desenvolvimento de algoritmos e técnicas de otimização para 
esse tipo de problema representa não apenas um avanço teórico, mas também uma oportunidade para soluções que possam transformar 
o setor de transporte e logística.

O objetivo principal deste trabalho é testar e comparar o desempenho das soluções para o problema de roteamento de veículos, 
utilizando diferentes técnicas de cálculo paralelo. Em particular, serão utilizadas plataformas de computação paralela como 
processadores gráficos (GPUs) e processadores de múltiplos núcleos (CPUs) para explorar e avaliar a eficiência dos métodos 
propostos. O uso de GPUs permite que cálculos intensivos sejam distribuídos entre milhares de núcleos de processamento, acelerando
operações complexas que seriam inviáveis com um único núcleo. Da mesma forma, a utilização de CPUs com múltiplos núcleos 
permite o processamento simultâneo de várias tarefas, o que é vantajoso para algoritmos que necessitam manipular grandes 
volumes de dados. A comparação entre essas tecnologias permitirá avaliar a escalabilidade e a adaptabilidade das soluções 
propostas, além de identificar as vantagens e os desafios associados a cada plataforma de cálculo paralelo. A análise dessas 
variáveis é essencial, pois permitirá uma compreensão mais profunda sobre o tipo de hardware que melhor se adapta a problemas 
complexos de roteamento de veículos.

A escolha do problema de roteamento de veículos (PRV) como objeto deste trabalho vai além de uma mera curiosidade pelo tema. 
Observando os ambientes em que já trabalhei, percebi que a falta de uma solução eficiente para esse tipo de problema é uma limitação 
para muitas empresas, tanto em termos de custo quanto de eficiência operacional. Este projeto representa uma oportunidade de colocar 
em prática os conhecimentos adquiridos durante e fora da graduação, especialmente nas áreas de computação de alto desempenho (HPC), 
algoritmos de otimização, programação paralela e utilização de ferramentas amplamente utilizadas no mercado. O projeto me permitirá 
aprofundar habilidades nessas áreas, com foco na aplicação prática e no uso de tecnologias de código aberto, sempre que possível. A 
escolha por ferramentas de código aberto reforça o compromisso com a acessibilidade e a democratização do conhecimento, permitindo 
que outros pesquisadores e profissionais possam beneficiar-se dos avanços realizados e aplicá-los em seus próprios contextos.

A experimentação com GPUs e CPUs também me permitirá uma compreensão mais aprofundada sobre a escolha de ferramentas adequadas 
para diferentes tipos de problemas de otimização, contribuindo para meu crescimento acadêmico e profissional.

\chapter{Objetivo}

O trabalho será desenvolvido com foco em uma solução de baixo custo e que possa ser replicada e testada por outros interessados, 
de modo a promover o avanço colaborativo e acessível na área de roteamento de veículos. Para isso, o projeto será iniciado no meu computador pessoal,
aproveitando-se de uma imagem em Docker, que facilitará a configuração do ambiente de desenvolvimento 
e permitirá a criação de um ambiente controlado e portátil. Esta abordagem tem a vantagem de garantir que o projeto possa ser executado 
em diferentes máquinas sem a necessidade de ajustes complexos, permitindo que outras pessoas reproduzam e validem os resultados de 
forma simples.

Após essa etapa inicial, o próximo passo será avaliar a viabilidade e explorar maneiras de transferir a execução do projeto para ambientes em cloud. 
Isso permitirá que o projeto se beneficie de recursos computacionais extras, através de clusters, que são especialmente úteis para o processamento de 
grandes volumes de dados e execução de operações intensivas, comuns em problemas de otimização de rotas. Ao testar o projeto nesses ambientes, 
será possível avaliar a escalabilidade das soluções desenvolvidas, verificando se o desempenho melhora proporcionalmente à medida que mais recursos 
de hardware são disponibilizados.

Além disso, caso viável, a adaptação para ambientes em nuvem permitirá a análise de diferentes configurações de hardware e software, possibilitando 
identificar quais fatores impactam o desempenho em cenários reais. Isso inclui investigar o impacto de diferentes arquiteturas de processadores. 
Com essas análises, o trabalho busca contribuir para a definição de melhores práticas e diretrizes para a execução de problemas de roteamento e implementações no setor de logística e transporte.

\bibliography{refs.bib}


\end{document}
